\documentclass[11pt, a4paper]{article}

% --- UNIVERSAL PREAMBLE BLOCK ---
\usepackage[top=2.5cm, bottom=2.5cm, left=2cm, right=2cm]{geometry}
\usepackage{fontspec}
\usepackage[english, bidi=basic, provide=*]{babel}
\babelfont{rm}{Noto Sans}

\usepackage{amsmath, amssymb, booktabs}
\usepackage{hyperref}

% --- Notation ---
\newcommand{\vswirl}{\mathbf{v}_{\!\boldsymbol{\circlearrowleft}}}
\newcommand{\rhof}{\rho_{\!f}}
\newcommand{\rhocore}{\rho_{\text{core}}}
\newcommand{\rc}{r_c}
\newcommand{\GammaSST}{\Gamma_0}
\newcommand{\Ltot}{L_{\text{tot}}}
\newcommand{\goldenratio}{\varphi}
\newcommand{\alphafs}{\alpha}

\title{\textbf{Gapped Excitations and Discrete Scale Invariance in a Vortex-Filament Toy Model:\\
A Thermodynamic Decomposition of an Excitation Hierarchy}}
\author{Omar Iskandarani\\ \small Independent Researcher, Groningen, The Netherlands\\
\small \href{mailto:info@omariskandarani.com}{info@omariskandarani.com} \quad ORCID: 0009-0006-1686-3961}
\date{March 2026}

\begin{document}
\maketitle

\noindent\textbf{Definitions:} $\goldenratio\equiv (1+\sqrt{5})/2$ and $\alphafs \equiv e^2/(4\pi\epsilon_0\hbar c)$.

\begin{abstract}
We study a vortex-filament \emph{toy model} in an idealized incompressible background and propose a thermodynamic bookkeeping of its excitation hierarchy.
Using the Abe--Okuyama quantum--thermodynamic isomorphism, we decompose filament energy changes into an ``adiabatic'' channel (geometry/topology preserved) and a ``heat'' channel (redistribution among internal Kelvin-like modes at fixed topology).
We further introduce a discrete-scale multiplier to parameterize a log-periodic excitation ladder consistent with discrete scale invariance.
As a numerical reference only, we compare a small set of toy-model excitation scales to several Particle Data Group mass scales; we do \emph{not} claim Standard Model particle identities.
\end{abstract}

\vspace{0.35cm}

%======================================================================
\section{Model scope and assumptions}
\label{sec:scope}
%======================================================================

We consider a thin vortex-filament configuration \(K\) embedded in an idealized inviscid, incompressible continuum.
The present paper is intentionally conservative in interpretation: all identifications with particle physics are treated as \emph{numerical reference scales} rather than literal ontology.
The model ingredients are:
(i) a relaxed geometric invariant \(\Ltot(K)\) obtained from a filament relaxation engine,
(ii) a phenomenological topological screening functional \(IM(K)\),
(iii) a gapped Kelvin-like internal excitation spectrum with gap \(\Delta_K\), and
(iv) a discrete-scale excitation index \(G\) controlling a log-periodic ladder.

%======================================================================
\section{Thermodynamic decomposition}
\label{sec:thermo}
%======================================================================

Following Abe--Okuyama, we use the identity
\begin{equation}
dE = \delta W + \delta Q,
\label{eq:firstlaw}
\end{equation}
with \(\delta W=\sum_k p_k\,dE_k\) and \(\delta Q=\sum_k E_k\,dp_k\), as a bookkeeping device separating changes of energy levels from changes of their occupation probabilities.

\subsection{Adiabatic baseline (geometry/topology preserved)}
We define a baseline energy (or mass-equivalent scale) for topology \(K\) as
\begin{equation}
M_{\mathrm{base}}(K) = \rho_0\,\chi(K)\,IM(K)\,\Ltot(K),
\label{eq:Mbase}
\end{equation}
where \(\rho_0\) is a single calibration constant, \(\chi(K)\in\{1,2\}\) is a parity factor (e.g.\ even/odd component convention), and
\begin{equation}
IM(K) \equiv b(K)^{-3/2}\,\goldenratio^{-g(K)}\,n(K)^{-1/\goldenratio}.
\label{eq:IM}
\end{equation}
Here \(b,g,n\) denote a chosen set of topological invariants (bridge/strand density, genus, component count). Equation \eqref{eq:Mbase} is a \emph{phenomenological kernel}; it is not derived from first-principles fluid energetics in this paper.

\subsection{Heat channel and a gapped Kelvin-like spectrum}
We assume internal Kelvin-like excitations are gapped by topological constraints, reconnection suppression, or torsional barriers, summarized by a topology-dependent excitation gap \(\Delta_K\).
In the low-temperature regime \(k_BT\ll \Delta_K\), excitation probabilities are exponentially suppressed,
\begin{equation}
p_{\mathrm{exc}} \propto e^{-\Delta_K/(k_BT)},
\qquad
C_V \propto e^{-\Delta_K/(k_BT)}.
\label{eq:gap_supp}
\end{equation}
This yields a frozen-response limit \(C_V\to 0\) at low temperature, analogous to standard gapped systems.

%======================================================================
\section{Discrete scale invariance and excitation index}
\label{sec:dsi}
%======================================================================

We parameterize a discrete excitation ladder by an integer index \(G\ge 0\) with a dimensionless multiplier
\begin{equation}
m(K,G) = M_{\mathrm{base}}(K)\,\Lambda(G),
\qquad
\Lambda(G) \equiv \left(\frac{4}{\alphafs}\right)^G.
\label{eq:mKG}
\end{equation}
We emphasize that \(\Lambda(G)\) is an \emph{ansatz} intended to compactly capture a large scale separation. Under discrete scale invariance, it is convenient to express \(\Lambda(G)\) in golden-ladder form,
\begin{equation}
\left(\frac{4}{\alphafs}\right)^G \approx \goldenratio^{\nu G},
\qquad
\nu \equiv \log_{\goldenratio}\!\left(\frac{4}{\alphafs}\right)\approx 13.1 \simeq 13,
\label{eq:nu}
\end{equation}
which provides a log-periodic interpretation of the excitation steps.

%======================================================================
\section{Illustrative examples (numerical reference scales)}
\label{sec:examples}
%======================================================================

We report a few illustrative cases using invariants and ropelength values returned by the relaxation engine.
The numerical values below are presented as \emph{reference scales} only.

\paragraph{Trefoil topology \(3_1\).}
Using
\[
b=2,\quad g=1,\quad n=1,\quad \chi=2,\quad \Ltot=46.450,
\]
one obtains example scales
\[
m(3_1,G=0)=0.4609~\mathrm{MeV},\qquad
m(3_1,G=2)=138.488~\mathrm{GeV}.
\]

\paragraph{Cinque-foil torus topology \(5_1\).}
Using
\[
b=2,\quad g=2,\quad n=1,\quad \chi=2,\quad \Ltot=57.410,
\]
one obtains
\[
m(5_1,G=1)=192.989~\mathrm{MeV},\qquad
m(5_1,G=2)=105.786~\mathrm{GeV}.
\]

\paragraph{Unknot topology \(1_1\).}
Using
\[
b=1,\quad g=0,\quad n=1,\quad \chi=2,\quad \Ltot=25.932,
\]
one obtains
\[
m(1_1,G=1)=645.508~\mathrm{MeV}.
\]

\subsection{Interpretation}
Within the toy model, different values of \(G\) represent distinct excitation levels at fixed topology \(K\). Any comparison to Standard Model masses is purely numerical and does not imply particle identity or field-theoretic equivalence.

%======================================================================
\section{Thermodynamic response and falsifiable signatures}
\label{sec:predictions}
%======================================================================

A generic prediction of the combined ``gap + log-periodic ladder'' structure is a crossover from exponentially suppressed response \eqref{eq:gap_supp} at low temperature to activated ladder behavior at higher temperature.
For a truncated golden ladder \(E_n=E_0\goldenratio^n\), the canonical heat capacity is
\begin{equation}
C_V = k_B \beta^2 \left[ \frac{1}{Z} \sum_{n=0}^N E_n^2 e^{-\beta E_n} - \left( \frac{1}{Z} \sum_{n=0}^N E_n e^{-\beta E_n} \right)^2 \right],
\qquad
Z=\sum_{n=0}^N e^{-\beta E_n}.
\label{eq:CV}
\end{equation}
Log-periodic features in response functions (e.g.\ oscillations vs \(\ln T\) or \(\ln E\)) are a falsifiable signature of discrete scale invariance.

%======================================================================
\section{Limitations}
\label{sec:limits}
%======================================================================

The present work is a phenomenological toy model. The kernel \eqref{eq:Mbase} and excitation multiplier \eqref{eq:mKG} are not derived from an Euler/Lagrangian action or a quantum field theory, and the numerical calibration \(\rho_0\) is empirical. Establishing a first-principles derivation would require an explicit Hamiltonian for the filament and a controlled computation of Kelvin-mode spectra and coupling.

%======================================================================
\section{Conclusion}
\label{sec:conclusion}
%======================================================================

We presented a thermodynamic decomposition of a vortex-filament excitation hierarchy with an explicit gap \(\Delta_K\) and a discrete-scale excitation index \(G\). The framework predicts frozen response at low temperature and log-periodic ladder activation at higher temperature. Comparisons to particle-physics mass scales were used only as numerical reference points.

\begin{thebibliography}{9}

\bibitem{AbeOkuyama2011}
Abe, S., \& Okuyama, S. (2011).
``Similarity between quantum mechanics and thermodynamics: Entropy, temperature, and Carnot cycle''.
\textit{Physical Review E}, 83, 021121.
DOI: 10.1103/PhysRevE.83.021121.

\bibitem{Sornette1998}
Sornette, D. (1998).
``Discrete scale invariance and complex dimensions''.
\textit{Physics Reports}, 297, 239--270.
DOI: 10.1016/S0370-1573(97)00076-8.

\bibitem{Saffman1992}
Saffman, P. G. (1992).
\textit{Vortex Dynamics}.
Cambridge University Press.
DOI: 10.1017/CBO9780511624063.

\bibitem{Batchelor1967}
Batchelor, G. K. (1967).
\textit{An Introduction to Fluid Dynamics}.
Cambridge University Press.

\bibitem{LandauLifshitz1987}
Landau, L. D., \& Lifshitz, E. M. (1987).
\textit{Fluid Mechanics}, 2nd ed. (Course of Theoretical Physics, Vol. 6).
Pergamon Press.

\end{thebibliography}

\end{document}
